\section*{अध्याय \ref{ch:g12:climate-change} --- जलवायु परिवर्तन: प्रेक्षण, कारण और प्रक्षेपण}

\begin{solution}{exo:g12:climate-change:1}
वह परिवर्तन जो वह वायुमंडल के शीर्ष पर जलवायु तंत्र में प्रवेश करने वाले ऊर्जा के शुद्ध प्रवाह में
उत्पन्न करता है, इससे पहले कि सतह का तापमान प्रतिक्रिया करे;
इसे \unit{W/m^2} में मापा जाता है। धनात्मक: ग्रह ऊर्जा प्राप्त करता है और गरम होता है;
ऋणात्मक: वह ऊर्जा खोता है और ठंडा होता है।
\end{solution}

\begin{solution}{exo:g12:climate-change:2}
मई में सबसे ऊँची, सितंबर में सबसे नीची। वसंत और गर्मियों में उत्तरी गोलार्ध की
वनस्पति, जहाँ अधिकांश स्थलीय पौधे उगते हैं, प्रकाशसंश्लेषण द्वारा उससे अधिक
कार्बन डाइऑक्साइड ग्रहण करती है जितनी पौधे और मृदाएँ छोड़ते हैं; पतझड़
और सर्दियों में अपघटन और श्वसन हावी रहते हैं और सांद्रता फिर से बढ़ती है।
\end{solution}

\begin{solution}{exo:g12:climate-change:3}
\emph{बहुत संभावित}: 90 से 100\,\% की प्रायिकता; \emph{संभावित}: 66 से
100\,\%। $[0.95 \text{ से } 1.20]$ एक बहुत संभावित परास है: 90\,\% संभावना कि
सही मान 0.95 और 1.20 के बीच है।
\end{solution}

\begin{solution}{exo:g12:climate-change:4}
91\,\%। जल गरम होने पर फैलता है: अधिक गरम महासागर अधिक आयतन घेरता है, इसलिए
उसकी सतह उठती है, चाहे कोई जल न जोड़ा जाए (स्थल की पिघलती बर्फ़ इसके ऊपर से
जल जोड़ती है)।
\end{solution}

\begin{solution}{exo:g12:climate-change:5}
जल-वाष्प: गरम वायु अधिक जल-वाष्प रखती है, जो एक ग्रीनहाउस गैस है, और इससे
तापन बढ़ता है। हिम और बर्फ़: पिघलने से गहरे रंग की सतहें खुल जाती हैं, जो
अधिक सूर्य-प्रकाश अवशोषित करती हैं। ऋणात्मक प्रतिक्रिया: अधिक गरम ग्रह अंतरिक्ष में अधिक अवरक्त
विकिरित करता है, जो उसे वापस संतुलन में लाता है।
\end{solution}

\begin{solution}{exo:g12:climate-change:6}
$427.35 - 315.98 = \qty{111.37}{ppm}$; $111.37/315.98 \approx 0.352$, यानी
लगभग 35\,\% की वृद्धि। 66 वर्षों में: $111.37/66 \approx
\qty{1.7}{ppm}$ प्रति वर्ष, औसतन।
\end{solution}

\begin{solution}{exo:g12:climate-change:7}
$F = 3.93 \times \ln(427.35/278)/\ln 2 = 3.93 \times 0.4300/0.6931 \approx
\qty{2.44}{W/m^2}$।
\end{solution}

\begin{solution}{exo:g12:climate-change:8}
$3.7 \times 30 = \qty{111}{mm}$, लगभग \qty{11}{cm}। $3.7/1.3 \approx 2.8$:
लगभग तीन गुना तेज़।
\end{solution}

\begin{solution}{exo:g12:climate-change:9}
$\Delta T = 2.72/1.16 \approx \qty{2.3}{\degreeCelsius}$। जलवायु
साम्य पर नहीं है: महासागर, जो अतिरिक्त ऊर्जा का लगभग 90\,\% ग्रहण करता है,
धीरे-धीरे गरम होता है, इसलिए सतह अभी \omterm{def:g12:climate-change:forcing}{प्रणोदन} के बराबर नहीं पहुँची है (तंत्र
अब भी लगभग \qty{0.8}{W/m^2} प्राप्त कर रहा है)।
\end{solution}

\begin{solution}{exo:g12:climate-change:10}
सूर्य और ज्वालामुखी \qty{1.06}{\degreeCelsius} के तापन में से अधिक से अधिक
\qty{0.1}{\degreeCelsius} की व्याख्या करते हैं: वे उसका कारण नहीं हो सकते। दूसरा
प्राकृतिक कारण जलवायु की आंतरिक परिवर्तनशीलता है (उसके प्राकृतिक
उतार-चढ़ाव, जैसे महासागर और वायुमंडल के बीच ऊष्मा का आदान-प्रदान); आकलन
उसका योगदान $-0.2$ से $+\qty{0.2}{\degreeCelsius}$ रखता है,
जो भी बहुत छोटा है।
\end{solution}

\begin{solution}{exo:g12:climate-change:11}
बहुत संभावित रूप से \qty{2}{\degreeCelsius} से नीचे: केवल SSP1-1.9 (बहुत संभावित परास
1.0 से 1.8)। बहुत संभावित रूप से ऊपर: SSP2-4.5 (2.1 से 3.5), SSP3-7.0 (2.8 से 4.6)
और SSP5-8.5 (3.3 से 5.7)। SSP1-2.6 (1.3 से 2.4) के लिए दोनों में से कोई भी बहुत संभावित नहीं है।
\end{solution}

\begin{solution}{exo:g12:climate-change:12}
एरोसॉल \omterm{def:g12:climate-change:forcing}{प्रणोदन} $-0.22 - 0.84 = \qty{-1.06}{W/m^2}$ है; उसके बिना
कुल $2.72 + 1.06 = \qty{3.78}{W/m^2}$ हो जाता है। साम्य तापन
$1.06/1.16 \approx \qty{0.9}{\degreeCelsius}$ बढ़ता है।
\end{solution}

\begin{solution}{exo:g12:climate-change:13}
चूँकि $\ln(2C) - \ln C = \ln 2$, $C$ कुछ भी हो,
\[
  F(556) - F(278) = \frac{3.93\,\ln 2}{\ln 2} = \qty{3.93}{W/m^2}, \qquad
  F(1112) - F(556) = \frac{3.93\,\ln 2}{\ln 2} = \qty{3.93}{W/m^2}:
\]
हर दोगुनेपन के लिए वही \omterm{def:g12:climate-change:forcing}{प्रणोदन}। \qty{1112}{ppm} पर कुल: $2 \times 3.93 =
\qty{7.86}{W/m^2}$।
\end{solution}

\begin{solution}{exo:g12:climate-change:14}
अवकलज है $\dfrac{\dd F}{\dd C} = \dfrac{F_{2\times}}{C \ln 2}$। \qty{280}{ppm} पर यह
$3.93/(280 \times 0.6931) \approx \qty{0.0202}{W/m^2}$ प्रति ppm;
\qty{420}{ppm} पर: $\approx \qty{0.0135}{W/m^2}$ प्रति ppm। सांद्रता बढ़ने के साथ हर अतिरिक्त ppm
कम \omterm{def:g12:climate-change:forcing}{प्रणोदन} जोड़ता है (280 की तुलना में 420 पर 1.5 गुना कम),
पर सांद्रता बढ़ती रहती है, इसलिए \omterm{def:g12:climate-change:forcing}{प्रणोदन} फिर भी बढ़ता है।
\end{solution}

\begin{solution}{exo:g12:climate-change:15}
\omterm{def:g12:climate-change:feedback}{प्रतिपुष्टियों} के बिना $3.93/3.22 \approx \qty{1.2}{\degreeCelsius}$;
उनके साथ $3.93/1.16 \approx \qty{3.4}{\degreeCelsius}$। प्रवर्धन
गुणक: $3.22/1.16 \approx 2.8$।
\end{solution}

\begin{solution}{pb:g12:climate-change:1}
\textbf{1.} $3.37 + 3.51 = \qty{6.88}{ppm}$, लगभग \qty{7}{ppm}।

\textbf{2.} वसंत से गर्मियों के अंत तक, उत्तरी
गोलार्ध के पौधे प्रकाशसंश्लेषण द्वारा वायु से कार्बन डाइऑक्साइड उससे तेज़ी से खींचते हैं
जितनी तेज़ी से श्वसन और अपघटन उसे लौटाते हैं: सांद्रता सितंबर तक गिरती है।
पतझड़ और सर्दियों में संतुलन उलट जाता है, और सांद्रता अपने मई के अधिकतम तक
फिर चढ़ जाती है।

\textbf{3.} दक्षिणी गोलार्ध में उत्तरी की तुलना में कहीं कम स्थल, और इसलिए कहीं कम
वनस्पति, है; उसकी अधिकांश सतह महासागर है।

\textbf{4.} ऊँचाई पर होने से, वह भूमि के निकट की प्रदूषित वायु की परत के ऊपर
है; महाद्वीपों और शहरों से दूर, जिस वायु का वह नमूना लेती है वह भली-भाँति मिश्रित है और
किसी स्थानीय स्रोत का नहीं, पूरे वायुमंडल का प्रतिनिधित्व करती है।

\textbf{5.} $(325.68 - 316.91)/10 \approx \qty{0.88}{ppm}$ प्रति वर्ष।

\textbf{6.} $(427.35 - 401.01)/10 \approx \qty{2.63}{ppm}$ प्रति वर्ष।

\textbf{7.} $2.634/0.877 \approx 3.0$: वृद्धि-दर तिगुनी हो गई है।

\textbf{8.} $(427.35 - 278)/278 \approx 0.537$: लगभग 54\,\%।

\textbf{9.} $427.35 - 300 \approx \qty{127}{ppm}$।

\textbf{10.} $F = 3.93 \times \ln(427.35/278)/\ln 2 \approx
\qty{2.44}{W/m^2}$।

\textbf{11.} $F(316.91) = 3.93 \times \ln(316.91/278)/\ln 2 \approx
\qty{0.74}{W/m^2}$; वृद्धि $2.44 - 0.74 \approx \qty{1.70}{W/m^2}$।

\textbf{12.} एक और \unit{W/m^2} सांद्रता को
$2^{1/3.93} \approx 1.193$ से गुणा करता है: $427.35 \times 1.193 \approx \qty{510}{ppm}$।

\textbf{13.} $(509.8 - 427.35)/2.634 \approx 31$ वर्ष।

\textbf{14.} $2 \times 278 = \qty{556}{ppm}$; \omterm{def:g12:climate-change:forcing}{प्रणोदन} $F_{2\times} =
\qty{3.93}{W/m^2}$।

\textbf{15.} संवेदनशीलता की परिभाषा से, \qty{3}{\degreeCelsius}
(संभावित रूप से 2.5 से \qty{4}{\degreeCelsius})।

\textbf{16.} $\lambda = 3.93/3 \approx \qty{1.31}{W.m^{-2}.\degreeCelsius^{-1}}$।
\omterm{def:g12:climate-change:feedback}{प्रतिपुष्टियों} का योग, \qty{1.16}{W.m^{-2}.\degreeCelsius^{-1}},
$3.93/1.16 \approx \qty{3.4}{\degreeCelsius}$ देता है, थोड़ा अधिक; आकलित
सर्वोत्तम अनुमान इस भौतिकी को प्रेक्षित तापन और अतीत की जलवायु के साथ
मिलाता है, और दोनों मान संभावित परास में पड़ते हैं।

\textbf{17.} महासागर अतिरिक्त ऊर्जा का अधिकांश भाग संचित करता है और धीरे-धीरे गरम होता है:
सतह अपने साम्य तापन तक दशकों से शताब्दियों के बाद ही पहुँचती है।

\textbf{18.} वह भविष्य के उत्सर्जनों की एक कहानी मान लेता है (पिछले दस वर्षों की
वृद्धि हमेशा बनी रहे): यह एक \omterm{def:g12:climate-change:scenario}{परिदृश्य} है। \omterm{def:g12:climate-change:scenario}{प्रक्षेपण} उत्सर्जनों के किसी \omterm{def:g12:climate-change:scenario}{परिदृश्य} पर
जलवायु और कार्बन चक्र की प्रतिक्रिया भौतिक मॉडलों से परिकलित करता है,
और परास और संभाव्यताएँ देता है।

\textbf{19.} 2025 के $(556 - 427.35)/2.634 \approx 48.8$ वर्ष बाद: लगभग
\textbf{2074}।
\end{solution}
